\documentclass{subfiles}
\begin{document}
    Consider the following scenario. 
        We are given a huge amount of data, which cannot be stored due to the volume.
        Our goal is that of processing such data to derive some statistics.

    Achiving such a goal isn't trivial; in fact, we have to change the usual model of 
        computation in favor of a streaming one.
        The idea is the following: we are given a sequence \(\sigma = \iprod{\sigma_{1}, \sigma_{2}, \ldots, \sigma_{m}}\)
        of tokens, each belonging to the same universe \(U\).
        The goal is to compute some function \(\Phi(\sigma)\) on the input stream.

    The simplest streaming model, the so-called \emph{vanilla} streaming model,
        assumes that the size of the universe \(U\), often treated as and integer universe
        (\(U = \set{0, \ldots, n - 1}\)), is known a priori, while the size of the stream 
        is unknown; therefore it requires that after each tokens is processed,
        we are able to return some value for \(\Phi(\sigma)\).

    As we shall see, is often useful to mantain a \emph{frequency vector} \(F_{\sigma}\),
        that keeps track of the number of occurrences of each token in the stream.
        In other words, \(F_{\sigma}[j]\) stores the number of occurrences of \(j\) in the stream.

    The analisis of a streaming algorithm differs a bit from the one usually done for algorithms.
        First, when it comes to storage, this is measured in terms of number of bits used 
        instead of by the number of elementary elements begin stored. Ideally,
        a streaming algorithm should use \(\ofO{\log(n + m)}\) bits of storage. 
        Second, the output of the algorithm often differs from the exact value of \(\Phi(\sigma)\).
        Hence another concern is the \emph{quality of the answer}.
        Furthermore, many streaming algorithms use \emph{randomization}, and may sometimes report an
        answer that does not lie within the desired error bounds. In such cases 
        the probability with which the algorithm gives a good approximation.

    \subsection{Frequent object}
    \subfile{../sub/4.1 - Frequent object}
\end{document}
